\documentclass[17pt,compress]{beamer}
\usepackage{beamerthemesplit}
\mode<presentation>
{
  \usetheme{Warsaw}
  \useoutertheme{infolines}
  \setbeamercovered{transparent}
  \setbeamertemplate{navigation symbols}{}
}
% Taken from Fernando's slides.
\usepackage{ae,aecompl}
\usepackage[scaled=.95]{helvet}

\usepackage[english]{babel}
\usepackage[latin1]{inputenc}
\usepackage[T1]{fontenc}

% change the alerted colour to LimeGreen
\definecolor{LimeGreen}{RGB}{50,205,50}
\setbeamercolor{structure}{fg=LimeGreen}
\author[FOSSEE]{}
\institute[IIT Bombay]{}
\date[]{}
% \setbeamercovered{transparent}

% theme split
\usepackage{verbatim}
\newenvironment{colorverbatim}[1][]%
{%
\color{blue}
\verbatim
}%
{%
\endverbatim
}%

\usepackage{mathpazo,courier,euler}
\usepackage{listings}
\lstset{language=sh,
    basicstyle=\ttfamily\bfseries,
  showstringspaces=false,
  keywordstyle=\color{black}\bfseries}

% logo
\logo{\includegraphics[height=1.30 cm]{3t-logo.pdf}}
\logo{\includegraphics[height=1.30 cm]{fossee-logo.pdf}

\hspace{7.5cm}
\includegraphics[scale=0.99]{fossee-logo.pdf}\\
\hspace{281pt}
\includegraphics[scale=0.80]{3t-logo.pdf}}

\begin{document}

\sffamily \bfseries
\title
[File Permissions and Ownership]
{File Permissions and Ownership}
\author
[FOSSEE]
{\small Talk to a Teacher\\{\color{blue}\url{http://spoken-tutorial.org}}\\\vspace{0.25cm}National Mission on Education
 through ICT\\{\color{blue}\url{ http://sakshat.ac.in}} \\ [0.77cm]
   Contributed by FOSSEE Team \\IIT Bombay  \\[0.3cm]
}

% slide 1
\begin{frame}
   \titlepage
\end{frame}

\begin{frame}
\frametitle{Objectives}
\label{sec-2}

At the end of this tutorial, you will be able to,
\begin{itemize}
\item Change file permissions
\item Change ownership of files
\end{itemize}
\end{frame}

\begin{frame}
\frametitle{Pre-requisites}
\label{sec-3}

Spoken tutorial on, 
\begin{itemize}
\item Getting started with Linux
\item Basic File Handling
\end{itemize}
\end{frame}

\begin{frame}[fragile]
  \frametitle{Symbolic modes}
  \begin{small}
    \begin{center}
      \begin{tabular}{lll}
        Ref.  &  Class   &  Description                                                      \\
        \hline
        u     &  user    &  the owner of the file                                            \\
        g     &  group   &  users who are members of the                                     \\
              &          &  file's group  \\
        o     &  others  &  users who are not the owner of the                           \\
              &          &  file or members of the group\\
        a     &  all     &  all three of the above; is the same                              \\
              &          &  as \emph{ugo}
      \end{tabular}
    \end{center}
   \end{small}
\end{frame}

\begin{frame}[fragile]
  \frametitle{Symbolic modes...}
  \begin{small}
    \begin{center}
      \begin{tabular}{ll}
        Operator  &  Description                                                                   \\
        \hline
         +         &  adds the specified modes to the               \\
                  &  specified classes                             \\
         -         &  removes the specified modes from the          \\
                  &  specified classes                             \\
         =         &  modes specified are to be made the            \\
                  &  exact modes for the specified classes  
      \end{tabular}
    \end{center}
   \end{small}
\end{frame}

\begin{frame}
  \frametitle{Symbolic modes...}
   \begin{small}
    \begin{center}
      \begin{tabular}{lll}
        Mode  &  Name     &  Description                                 \\
        \hline
        r     &  read     &  read a file or list a directory's   \\
              &           &  contents \\
        w     &  write    &  write to a file or directory                \\\vspace{2pt}
        x     &  execute  &  execute a file or recurse a   \\
              &           &  directory tree
      \end{tabular}
    \end{center}
  \end{small}
\end{frame}


\begin{frame}
  \frametitle{Exercise}
  \begin{itemize}
  \item Change the permissions of a directory including,
  \begin{itemize}
  \item All it's sub-directories
  \item All files within the directory
  \end{itemize} 
  \end{itemize}
\end{frame}

\begin{frame}
  \frametitle{Solution}
  \begin{itemize}
  \item \$ chmod go-r -R <directory name>/
  \end{itemize}
\end{frame}

\begin{frame}
\frametitle{Summary}

  In this tutorial, we have learnt to,

\begin{itemize}
\item Chane the permissions of files using the ``chmod'' command
\item Use the ``chown'' command to change the ownership of files
\end{itemize}
\end{frame}

\begin{frame}[fragile]
\frametitle{Evaluation}
\label{sec-9}


\begin{enumerate}
\item For a given file, change mode to r, w, x for all (user, group, others)
\vspace{10pt}
\item What changes, on specifying only an owner in the "chown" command?
\begin{itemize}
\item Only the owner of the file
\item The group ownership of the file
\item Neither the owner nor the group
\end{itemize}
\end{enumerate}
\end{frame}
\begin{frame}
\frametitle{Solutions}
\label{sec-10}


\begin{enumerate}
\item \$ chmod ugo+rwx <filename>
\vspace{15pt}
\item Only the owner of the file
\end{enumerate}
\end{frame}

\begin{frame}
\frametitle{SDES \& FOSSEE}
\begin{center}
\begin{itemize}
\item \small{SDES}\\
\small{\color{LimeGreen}Software Development techniques for Engineers and Scientists} \\
\scriptsize An initiative by FOSSEE. \\
\vspace{3pt}
\scriptsize For more information on SDES, please visit {\color{blue}\url{http://fossee.in/sdes}}\\
\vspace{10pt}
\item \small{FOSSEE}\\
\small {\color{LimeGreen}Free and Open-source Software for \\Science and Engineering Education} \\
\scriptsize Based at IIT Bombay, Funded by MHRD.\\
\vspace{3pt}
\scriptsize Part of National Mission on Education through ICT \\(NME-ICT) \\
\end{itemize}
\end{center}
\end{frame}

\begin{frame}
\frametitle{About the Spoken Tutorial Project}
\begin{itemize}
\item Watch the video available at {\color{blue}\url{http://spoken-tutorial.org /What\_is\_a\_Spoken\_Tutorial}} 
\item It summarises the Spoken Tutorial project 
\item If you do not have good bandwidth, you can download and watch it
\end{itemize}
\end{frame}

\begin{frame}
\frametitle{Spoken Tutorial Workshops}The Spoken Tutorial Project Team 
\begin{itemize}
\item Conducts workshops using spoken tutorials 
\item Gives certificates to those who pass an online test 
\item For more details, please write to \\ \hspace {0.5cm}{\color{blue}contact@spoken-tutorial.org}
\end{itemize}
\end{frame}

\begin{frame}
\frametitle{Acknowledgements}
\begin{itemize}
\item Spoken Tutorial Project is a part of the Talk to a Teacher  project 
\item It is supported by the National Mission on Education through  ICT, MHRD, Government of India 
\item More information on this Mission is available at: \\{\color{blue}\url{http://spoken-tutorial.org/NMEICT-Intro}}
\end{itemize}
\end{frame}

\begin{frame}

  \begin{block}{}
  \begin{center}
  {\Large THANK YOU!} 
  \end{center}
  \end{block}
\begin{block}{}
  \begin{center}
    For more Information, visit our website\\
    {\color{blue}\url{http://fossee.in/}}
  \end{center}  
  \end{block}
\end{frame}

\end{document}
